\begin{summarybox}
	\textbf{\textcolor{CSummary}{一句话核心}}
	\textbf{极线方程由曲线方程替换得到：平方变乘积，一次变平均；配极原理刻画了极点与极线的对称对应关系。}

	\medskip
	\textbf{\textcolor{CSummary}{使用条件}}
	\begin{itemize}[leftmargin=*, itemsep=0.5em, topsep=0.4em]
		\item \textbf{条件 1（标准方程）：} 曲线必须是椭圆 $\dfrac{x^2}{a^2}+\dfrac{y^2}{b^2}=1$、双曲线 $\dfrac{x^2}{a^2}-\dfrac{y^2}{b^2}=1$ 或抛物线 $y^2=2px$ 的标准形式，且非退化。
		\item \textbf{条件 2（已知点）：} 需要已知点 $P(x_0,y_0)$ 的坐标，代入公式时注意对应项。
	\end{itemize}

	\medskip
	\textbf{\textcolor{CSummary}{关键提醒}}
	\begin{itemize}[leftmargin=*, itemsep=0.5em, topsep=0.4em]
		\item \textbf{易错点（切线切点弦）：} 点在曲线上时极线是切线，点在曲线外时极线是切点弦，不可混淆。
		\item \textbf{检查项（符号系数）：} 双曲线极线方程是减号，抛物线极线方程中 $p$ 的系数是 $p$ 而非 $2p$。
	\end{itemize}

\end{summarybox}
